% !TEX program = xelatex
% ============================================================
%  Arabic Lab Report Template — Nibras Center
%  Compile with: xelatex main.tex (twice)
% ============================================================
%  Features:
%    - A4, 11pt, article class
%    - Title block: experiment name, course, date, instructor
%    - Sections: objective, theory, equipment, procedure,
%      results, analysis, conclusion
%    - Data tables and figures
% ============================================================

\documentclass[11pt,a4paper]{article}

% ---- Geometry ----
\usepackage[a4paper,margin=2.5cm,top=2.5cm,bottom=2.5cm,headheight=15pt]{geometry}

% ---- Core packages (order matters for bidi) ----
\usepackage{url}
\usepackage{amsmath}
\usepackage{amssymb}
\usepackage{graphicx}
\usepackage{array}
\usepackage{booktabs}
\usepackage{tabularx}
\usepackage{multirow}
\usepackage{enumitem}
\usepackage{float}
\usepackage{caption}
\usepackage{titlesec}
\usepackage{fancyhdr}
\usepackage{setspace}
\usepackage{siunitx}

% ---- RTL / Arabic language support ----
\usepackage{bidi}
\usepackage[no-math]{fontspec}
\usepackage[quiet]{polyglossia}

% ---- Fonts ----
\setmainfont[Scale=1]{Amiri-Regular.ttf}
\newfontfamily\arabicfont[Script=Arabic,Scale=0.95]{Amiri-Regular.ttf}
\newfontfamily\englishfont[Scale=0.95]{TeX Gyre Termes}
\setmonofont{Amiri-Regular.ttf}

% ---- Languages ----
\setdefaultlanguage[calendar=gregorian,hijricorrection=1]{arabic}
\setotherlanguage[variant=british]{english}

% ---- Hyperref ----
\usepackage[breaklinks,hidelinks]{hyperref}
\hypersetup{
  pdftitle={تقرير مختبر},
  pdfauthor={اسم الطالب}
}

% ---- Captions in Arabic ----
\addto\captionsarabic{%
  \renewcommand{\contentsname}{الفهرس}%
  \renewcommand{\figurename}{الشكل}%
  \renewcommand{\tablename}{الجدول}%
  \renewcommand{\refname}{المراجع}%
}

% ---- Section formatting ----
\titleformat{\section}{\Large\bfseries}{\thesection}{0.7em}{}
\titleformat{\subsection}{\large\bfseries}{\thesubsection}{0.6em}{}
\titlespacing*{\section}{0pt}{1.4ex plus 0.4ex}{0.9ex plus 0.3ex}
\titlespacing*{\subsection}{0pt}{1.0ex plus 0.3ex}{0.7ex plus 0.2ex}

% ---- Headers / footers ----
\pagestyle{fancy}
\fancyhf{}
\fancyhead[R]{\small\itshape تقرير مختبر: عنوان التجربة}
\fancyhead[L]{\small اسم الطالب}
\fancyfoot[C]{\small\thepage}
\renewcommand{\headrulewidth}{0.3pt}
\renewcommand{\footrulewidth}{0pt}

% ---- List spacing ----
\setlist[itemize]{topsep=0.3em, itemsep=0.15em, parsep=0pt}
\setlist[enumerate]{topsep=0.3em, itemsep=0.15em, parsep=0pt}

% ---- Table column types ----
\newcolumntype{L}[1]{>{\raggedright\arraybackslash}p{#1}}
\newcolumntype{C}[1]{>{\centering\arraybackslash}p{#1}}

% ---- Line spacing ----
\setstretch{1.5}

\begin{document}

% ============================================================
% TITLE BLOCK
% ============================================================
\begin{center}
{\LARGE\bfseries تقرير تجربة مختبرية}\\[1em]
{\Large عنوان التجربة}\\[1.5em]

\begin{tabular}{rl}
\textbf{اسم الطالب:} & --- \\
\textbf{الرقم الجامعي:} & --- \\
\textbf{المقرر:} & --- \\
\textbf{اسم المدرّس:} & --- \\
\textbf{تاريخ التجربة:} & --- \\
\textbf{تاريخ التسليم:} & --- \\
\end{tabular}
\end{center}

\vspace{1.5em}

% ============================================================
% 1. OBJECTIVE
% ============================================================
\section{الهدف من التجربة}
\label{sec:objective}

% TODO: اكتب هدف التجربة هنا في 2-3 جمل.
هذا نص تجريبي لهدف التجربة. يصف الهدف الرئيسي من إجراء التجربة والمهارات المستهدفة.

% ============================================================
% 2. THEORY
% ============================================================
\section{الأساس النظري}
\label{sec:theory}

% TODO: اكتب الأساس النظري للتجربة هنا.
% اذكر المعادلات الأساسية مع أرقامها.
هذا نص تجريبي للأساس النظري. يشرح المبادئ العلمية التي تقوم عليها التجربة.

\begin{equation}
\label{eq:main}
F = m \cdot a
\end{equation}
حيث $F$ القوة بوحدة \LR{Newton}، $m$ الكتلة بوحدة \LR{kg}، $a$ التسارع بوحدة \LR{m/s$^2$}.

% ============================================================
% 3. EQUIPMENT
% ============================================================
\section{الأدوات والمعدات}
\label{sec:equipment}

% TODO: اذكر الأدوات والمعدات المستخدمة.
\begin{itemize}
    \item جهاز قياس
    \item أسلاك توصيل
    \item مصدر طاقة
    \item أدوات قياس أخرى
\end{itemize}

% ============================================================
% 4. PROCEDURE
% ============================================================
\section{خطوات العمل}
\label{sec:procedure}

% TODO: اكتب خطوات العمل بالترتيب.
\begin{enumerate}
    \item الخطوة الأولى: وصف ما تم عمله.
    \item الخطوة الثانية: وصف ما تم عمله.
    \item الخطوة الثالثة: وصف ما تم عمله.
    \item الخطوة الرابعة: وصف ما تم عمله.
\end{enumerate}

% ============================================================
% 5. RESULTS
% ============================================================
\section{النتائج}
\label{sec:results}

% TODO: اعرض النتائج في جداول وأشكال.

\begin{table}[H]
\centering
\caption{بيانات القياس}
\begin{tabular}{|c|c|c|c|}
\hline
\textbf{القراءة} & \textbf{القيمة 1} & \textbf{القيمة 2} & \textbf{المتوسط} \\
\hline
1 & 10.0 & 10.2 & 10.1 \\
2 & 20.0 & 19.8 & 19.9 \\
3 & 30.0 & 30.1 & 30.05 \\
4 & 40.0 & 39.9 & 39.95 \\
5 & 50.0 & 50.2 & 50.1 \\
\hline
\end{tabular}
\end{table}

% ---- Example figure ----
% \begin{figure}[H]
% \centering
% \includegraphics[width=0.7\textwidth]{plot.png}
% \caption{منحنى العلاقة بين المتغيرين}
% \label{fig:plot}
% \end{figure}

% ============================================================
% 6. ANALYSIS
% ============================================================
\section{التحليل والمناقشة}
\label{sec:analysis}

% TODO: حلل النتائج وقارنها بالقيم النظرية.
هذا نص تجريبي للتحليل. يحلل النتائج التجريبية ويقارنها بالقيم النظرية المتوقعة.

\subsection{حساب نسبة الخطأ}
% TODO: احسب نسبة الخطأ بين القيم التجريبية والنظرية.
نسبة الخطأ المحسوبة: \LR{2.5\%}.

\subsection{مصادر الخطأ}
% TODO: اذكر مصادر الخطأ المحتملة.
\begin{itemize}
    \item خطأ في قراءة الأجهزة.
    \item خطأ في ضبط الصفر.
    \item تأثير العوامل البيئية.
\end{itemize}

% ============================================================
% 7. CONCLUSION
% ============================================================
\section{الاستنتاج}
\label{sec:conclusion}

% TODO: اكتب الاستنتاج النهائي.
هذا نص تجريبي للاستنتاج. يلخص النتائج الرئيسية ويؤكد ما إذا كانت الأهداف قد تحققت.

% ============================================================
% REFERENCES (optional)
% ============================================================
\section*{المراجع}
\addcontentsline{toc}{section}{المراجع}

\begin{thebibliography}{9}
\bibitem{ref1} Author, A. (2024). \emph{Title}. Journal, 1(1), 1--10.
\end{thebibliography}

\end{document}
